% GENERATED FILE. Edit canonical lessons, then run npm run generate:books.
% canonical-source-hash: fnv1a64:ba665fdbafa4396d
% canonical-lessons: TA-C343-palankatai, TA-C343-puranam, TA-C343-kattukkatai, TA-C343-punai, TA-C343-palamoli

\chapter{Legends}
\label{ch:ta-legends}
\hlchaptermodality{343}

\begin{chapteropening}
\textbf{By the end of this chapter:} \emph{I can talk about legends, myths and proverbs.}

Complete the last lesson of chapter 343: I can talk about legends, myths and proverbs.
\end{chapteropening}

\section[\texorpdfstring{\ta{பழங்கதை} (paḻaṅkatai)}{பழங்கதை (paḻaṅkatai)}]{\ta{பழங்கதை} (paḻaṅkatai) — a legend; an old tale}

\label{lesson:TA-C343-palankatai}



\begin{quote}
Before the new one: say the Tamil for an age, then the Tamil for archaeology.
\end{quote}

\subsection*{You'll want to know: \ta{பழங்கதை}}
\textbf{\ta{பழங்கதை}} — \emph{paḻaṅkatai} — \textquotedblleft{}a legend; an old tale\textquotedblright{}.

\begin{quote}
\ta{பாட்டி} \ta{ஒரு} \ta{பழங்கதை} \ta{சொன்னார்}.
\end{quote}

\begin{quote}
\emph{pāṭṭi oru paḻaṅkatai coṉṉār.}
\end{quote}

\begin{quote}
Grandmother told an old legend.
\end{quote}

\ta{பழம்}, old + \ta{கதை}, story.

\begin{cousinweb}
\textbf{\ta{பழங்கதை}} (\emph{paḻaṅkatai}) joins native \ta{பழ}, \textquotedblleft{}old\textquotedblright{} (DEDR 3999), with \ta{கதை}, \textquotedblleft{}story\textquotedblright{}, from Sanskrit \emph{kathā} (CDIAL 2704).
\end{cousinweb}

\begin{grammarlens}[title={What you've built}]
One more word for telling what happened.
\end{grammarlens}

\subsection*{Your turn}
\begin{itemize}
\raggedright
  \item \emph{Say it:} \emph{paḻaṅkatai}
  \item \emph{Say it:} \emph{paḻaṅkatai}, once more
  \item \emph{Say it:} \emph{pāṭṭi oru paḻaṅkatai coṉṉār.} — Grandmother told an old legend.
  \item \emph{Recall:} say the Tamil for an age, then the Tamil for archaeology, then say \emph{paḻaṅkatai} again
\end{itemize}

\begin{tcolorbox}[breakable,colback=teal!4,colframe=teal!35!black,title={Before you move on}]
What does \ta{பழங்கதை} mean? (\textquotedblleft{}A legend; an old tale\textquotedblright{}.) Say it once more.
\end{tcolorbox}
\section[\texorpdfstring{\ta{புராணம்} (purāṇam)}{புராணம் (purāṇam)}]{\ta{புராணம்} (purāṇam) — a myth; a sacred legend}

\label{lesson:TA-C343-puranam}



\begin{quote}
Before the new one: say the Tamil for archaeology, then the Tamil for a legend.
\end{quote}

\subsection*{You'll want to know: \ta{புராணம்}}
\textbf{\ta{புராணம்}} — \emph{purāṇam} — \textquotedblleft{}a myth; a sacred legend\textquotedblright{}.

\begin{quote}
\ta{இந்த} \ta{ஊருக்கு} \ta{ஒரு} \ta{பழைய} \ta{புராணம்} \ta{இருக்கிறது}.
\end{quote}

\begin{quote}
\emph{inta ūrukku oru paḻaiya purāṇam irukkiṟatu.}
\end{quote}

\begin{quote}
This town has an old myth about it.
\end{quote}

Plural \ta{புராணங்கள்}.

\begin{cousinweb}
\textbf{\ta{புராணம்}} (\emph{purāṇam}) is borrowed from Sanskrit \emph{purāṇa}, \textquotedblleft{}ancient\textquotedblright{} (CDIAL 8283): an old tale.
\end{cousinweb}

\begin{grammarlens}[title={What you've built}]
One more word for telling what happened.
\end{grammarlens}

\subsection*{Your turn}
\begin{itemize}
\raggedright
  \item \emph{Say it:} \emph{purāṇam}
  \item \emph{Say it:} \emph{purāṇam}, once more
  \item \emph{Say it:} \emph{inta ūrukku oru paḻaiya purāṇam irukkiṟatu.} — This town has an old myth about it.
  \item \emph{Recall:} say the Tamil for archaeology, then the Tamil for a legend, then say \emph{purāṇam} again
\end{itemize}

\begin{tcolorbox}[breakable,colback=teal!4,colframe=teal!35!black,title={Before you move on}]
What does \ta{புராணம்} mean? (\textquotedblleft{}A myth; a sacred legend\textquotedblright{}.) Say it once more.
\end{tcolorbox}
\section[\texorpdfstring{\ta{கட்டுக்கதை} (kaṭṭukkatai)}{கட்டுக்கதை (kaṭṭukkatai)}]{\ta{கட்டுக்கதை} (kaṭṭukkatai) — a made-up story; a myth}

\label{lesson:TA-C343-kattukkatai}



\begin{quote}
Before the new one: say the Tamil for a legend, then the Tamil for a myth.
\end{quote}

\subsection*{You'll want to know: \ta{கட்டுக்கதை}}
\textbf{\ta{கட்டுக்கதை}} — \emph{kaṭṭukkatai} — \textquotedblleft{}a made-up story; a myth\textquotedblright{}.

\begin{quote}
\ta{அவன்} \ta{சொன்னது} \ta{ஒரு} \ta{கட்டுக்கதை}.
\end{quote}

\begin{quote}
\emph{avaṉ coṉṉatu oru kaṭṭukkatai.}
\end{quote}

\begin{quote}
What he said was a made-up story.
\end{quote}

\ta{கட்டு}, to build + \ta{கதை}: a story put together.

\begin{cousinweb}
\textbf{\ta{கட்டுக்கதை}} (\emph{kaṭṭukkatai}) joins native \ta{கட்டு}, \textquotedblleft{}to tie, build\textquotedblright{} (DEDR 1147), with \ta{கதை}, \textquotedblleft{}story\textquotedblright{}, from Sanskrit \emph{kathā} (CDIAL 2704).
\end{cousinweb}

\begin{grammarlens}[title={What you've built}]
One more word for telling what happened.
\end{grammarlens}

\subsection*{Your turn}
\begin{itemize}
\raggedright
  \item \emph{Say it:} \emph{kaṭṭukkatai}
  \item \emph{Say it:} \emph{kaṭṭukkatai}, once more
  \item \emph{Say it:} \emph{avaṉ coṉṉatu oru kaṭṭukkatai.} — What he said was a made-up story.
  \item \emph{Recall:} say the Tamil for a legend, then the Tamil for a myth, then say \emph{kaṭṭukkatai} again
\end{itemize}

\begin{tcolorbox}[breakable,colback=teal!4,colframe=teal!35!black,title={Before you move on}]
What does \ta{கட்டுக்கதை} mean? (\textquotedblleft{}A made-up story; a myth\textquotedblright{}.) Say it once more.
\end{tcolorbox}
\section[\texorpdfstring{\ta{புனை} (puṉai)}{புனை (puṉai)}]{\ta{புனை} (puṉai) — to make up; to compose}

\label{lesson:TA-C343-punai}



\begin{quote}
Before the new one: say the Tamil for a myth, then the Tamil for a made-up story.
\end{quote}

\subsection*{You'll want to know: \ta{புனை}}
\textbf{\ta{புனை}} — \emph{puṉai} — \textquotedblleft{}to make up; to compose\textquotedblright{}.

\begin{quote}
\ta{அவன்} \ta{ஒரு} \ta{பொய்யான} \ta{காரணத்தைப்} \ta{புனைந்தான்}.
\end{quote}

\begin{quote}
\emph{avaṉ oru poyyāṉa kāraṇattaip puṉaintāṉ.}
\end{quote}

\begin{quote}
He made up a false reason.
\end{quote}

Past (I …): \ta{புனைந்தேன்}, \emph{puṉaintēṉ}.

Also to compose a poem; \ta{புனைபெயர்}, a pen name.

\begin{cousinweb}
\textbf{\ta{புனை}} (\emph{puṉai}) is native Dravidian (DEDR 4340): to dress, adorn, make, form.
\end{cousinweb}

\begin{grammarlens}[title={What you've built}]
One more word for telling what happened.
\end{grammarlens}

\subsection*{Your turn}
\begin{itemize}
\raggedright
  \item \emph{Say it:} \emph{puṉai}
  \item \emph{Say it:} \emph{puṉai}, once more
  \item \emph{Say it:} \emph{avaṉ oru poyyāṉa kāraṇattaip puṉaintāṉ.} — He made up a false reason.
  \item \emph{Recall:} say the Tamil for a myth, then the Tamil for a made-up story, then say \emph{puṉai} again
\end{itemize}

\begin{tcolorbox}[breakable,colback=teal!4,colframe=teal!35!black,title={Before you move on}]
What does \ta{புனை} mean? (\textquotedblleft{}To make up; to compose\textquotedblright{}.) Say it once more.
\end{tcolorbox}
\section[\texorpdfstring{\ta{பழமொழி} (paḻamoḻi)}{பழமொழி (paḻamoḻi)}]{\ta{பழமொழி} (paḻamoḻi) — a proverb}

\label{lesson:TA-C343-palamoli}



\begin{quote}
Before the new one: say the Tamil for a made-up story, then the Tamil for to make up.
\end{quote}

\subsection*{You'll want to know: \ta{பழமொழி}}
\textbf{\ta{பழமொழி}} — \emph{paḻamoḻi} — \textquotedblleft{}a proverb\textquotedblright{}.

\begin{quote}
\ta{அவர்} \ta{பேசும்போது} \ta{நிறையப்} \ta{பழமொழிகளைச்} \ta{சொல்வார்}.
\end{quote}

\begin{quote}
\emph{avar pēcumpōtu niṟaiyap paḻamoḻikaḷaic colvār.}
\end{quote}

\begin{quote}
When he talks he uses lots of proverbs.
\end{quote}

Plural \ta{பழமொழிகள்}.

\begin{cousinweb}
\textbf{\ta{பழமொழி}} (\emph{paḻamoḻi}) is a Tamil compound built from \ta{பழ}, \textquotedblleft{}old\textquotedblright{} (DEDR 3999), and \ta{மொழி}, \textquotedblleft{}word, saying\textquotedblright{} (DEDR 4989).
\end{cousinweb}

\begin{grammarlens}[title={What you've built}]
One more word for telling what happened.
\end{grammarlens}

\subsection*{Your turn}
\begin{itemize}
\raggedright
  \item \emph{Say it:} \emph{paḻamoḻi}
  \item \emph{Say it:} \emph{paḻamoḻi}, once more
  \item \emph{Say it:} \emph{avar pēcumpōtu niṟaiyap paḻamoḻikaḷaic colvār.} — When he talks he uses lots of proverbs.
  \item \emph{Recall:} say the Tamil for a made-up story, then the Tamil for to make up, then say \emph{paḻamoḻi} again
\end{itemize}

\begin{tcolorbox}[breakable,colback=teal!4,colframe=teal!35!black,title={Before you move on}]
What does \ta{பழமொழி} mean? (\textquotedblleft{}A proverb\textquotedblright{}.) Say it once more.
\end{tcolorbox}
